% Informe técnico del Proyecto Orion · Optimización 2026-II · Universidad de los Llanos
% Compilar con pdflatex (Overleaf: plantilla IEEEtran de conferencia). Máximo 8 páginas.
\documentclass[conference]{IEEEtran}
\usepackage[utf8]{inputenc}
\usepackage[T1]{fontenc}
\usepackage[spanish,es-noshorthands]{babel}
\usepackage{amsmath,amssymb}
\usepackage{graphicx}
\usepackage{booktabs}
\usepackage{cite}
\usepackage{url}

\begin{document}

\title{Optimización de la red logística de EcoResiduos del Llano: modelo exacto, algoritmo genético y criticidad estructural}

% TODO(equipo): completar nombres completos y correos institucionales.
\author{\IEEEauthorblockN{David F. Beltrán Tabares (Líder)}
\IEEEauthorblockN{Julián [Apellidos] (Modelado y Optimización)}
\IEEEauthorblockN{Martín [Apellidos] (Metaheurísticas)}
\IEEEauthorblockN{Juan David [Apellidos] (Software e Interfaz)}
\IEEEauthorblockA{Grupo Orion · Optimización 2026-II\\
Facultad de Ciencias Básicas e Ingeniería, Universidad de los Llanos\\
Villavicencio, Colombia}}

\maketitle

\begin{abstract}
EcoResiduos del Llano recoge residuos agroindustriales (palma, arroz, ganadería y forestal) en el
departamento del Meta y los lleva a plantas de aprovechamiento. La asignación empírica de rutas
genera sobrecostos. Este trabajo modela la red como un grafo dirigido y resuelve el problema de flujo
de costo mínimo con dos métodos: un modelo de programación lineal (PuLP + CBC) y un algoritmo genético
(AG) con reparación de factibilidad. El sistema permite editar la red en línea, analizar escenarios
de sensibilidad (combustible, oferta y cuellos de botella), obtener precios sombra y rangos de
estabilidad, y comparar ambos métodos en costo, tiempo y brecha. Además, se incorporan centros de
acopio como nodos de paso y se adapta el Índice de Criticidad Estructural (ICE) para identificar las
plantas cuya falla más afecta a la red. En el caso base, el óptimo es de 64\,691\,000~COP por semana,
el AG lo alcanza con brecha nula, abrir dos acopios ahorra 4{,}5\,\% y el ICE ordena las plantas igual
que el impacto observado al simular su falla ($\rho = 0{,}80$ en el caso base y $\rho = 1{,}00$ en el
caso real del Meta).
\end{abstract}

\begin{IEEEkeywords}
flujo de costo mínimo, programación lineal, algoritmo genético, análisis de sensibilidad,
criticidad estructural, logística de residuos.
\end{IEEEkeywords}

\section{Introducción y problemática}
La agroindustria del Meta genera cada semana cientos de toneladas de residuos aprovechables: raquis y
fibra de palma, cascarilla de arroz, estiércol bovino y residuos forestales. EcoResiduos del Llano los
recoge en las empresas productoras y los lleva a cuatro plantas de procesamiento (Villavicencio,
Restrepo, Granada y Puerto Gaitán), cada una con capacidad y costo de operación propios. Hoy la
asignación empresa–planta se hace de forma empírica, lo que satura unas plantas y subutiliza otras.

El objetivo es un sistema que determine la distribución de costo mínimo, que se pueda ampliar con
nuevas empresas y sedes sin tocar el código y que permita evaluar en vivo qué pasa si sube el
combustible, crece la oferta o falla una planta. El trabajo cubre las cuatro fases del enunciado:
formulación (Fase~1), resolución exacta (Fase~2), algoritmo genético (Fase~3) e interfaz con análisis
de sensibilidad y comparación (Fase~4).

\section{Equipo y roles}
\begin{table}[h]
\centering
\caption{Roles del Grupo Orion}
\begin{tabular}{@{}ll@{}}
\toprule
Rol del enunciado & Responsabilidad principal \\
\midrule
Líder & Integración, ICE, informe y sustentación \\
Modelado y Optimización & Formulación, modelo exacto, duales y rangos \\
Metaheurísticas & Diseño y calibración del AG \\
Software e Interfaz & API, interfaz web, mapa y despliegue \\
\bottomrule
\end{tabular}
\end{table}

\section{Formulación matemática}
\subsection{Grafo de la red}
La red es un grafo dirigido $G=(V,E)$ con $V = O \cup T \cup P$: fuentes (empresas) $O$, centros de
acopio $T$ (nodos de paso) y plantas $P$. Los arcos son
$E \subseteq (O\times P) \cup (O\times T) \cup (T\times P)$; solo existen los arcos habilitados.

\subsection{Parámetros y variables}
$a_i$: oferta semanal de la empresa $i$ (t). $b_j$: capacidad de la planta $j$ (t). $u_t$: capacidad
del acopio $t$. $d_e$: distancia vial del arco $e$ (km). $k_e$: costo de flete por t·km
(350~COP por defecto; 200 en el flete consolidado acopio–planta). $o_j$: costo de operación por
tonelada en la planta $j$. $h_t$: costo de manejo en el acopio $t$. El costo unitario del arco es
$c_e = k_e d_e + o_{\text{destino}}$ o $c_e = k_e d_e + h_t$ si llega a un acopio. La variable de
decisión es $x_e \ge 0$, toneladas por el arco $e$.

\subsection{Modelo}
\begin{align}
\min\; Z &= \sum_{e\in E} c_e\, x_e \\
\text{s.a.}\; \sum_{e\in\delta^+(i)} x_e &= a_i && \forall i\in O \\
\sum_{e\in\delta^-(t)} x_e &= \sum_{e\in\delta^+(t)} x_e && \forall t\in T \\
\sum_{e\in\delta^-(t)} x_e &\le u_t && \forall t\in T \\
\sum_{e\in\delta^-(j)} x_e &\le b_j && \forall j\in P \\
x_e &\ge 0 && \forall e\in E
\end{align}
donde $\delta^+(v)$ y $\delta^-(v)$ son los arcos que salen y entran de $v$. Si
$\sum_i a_i > \sum_j b_j$, se agrega una planta ficticia con capacidad igual al exceso y una
penalización de $10^6$~COP/t, de modo que el modelo siempre es factible y las toneladas sin planta
quedan a la vista. Sin acopios ($T=\emptyset$), el modelo se reduce exactamente al problema de
transporte clásico, que es la prueba de regresión del sistema.

\section{Diseño del algoritmo genético}
\textbf{Cromosoma.} Una matriz empresa $\times$ camino, donde un camino es $i\to j$ o
$i\to t\to j$, con costo igual a la suma de sus arcos. Sin acopios, la matriz es la de transporte
$m\times n$. Cada fila reparte la oferta de una empresa.

\textbf{Población inicial.} El 90\,\% de los individuos reparte la oferta con una distribución de
Dirichlet sobre los caminos habilitados; el 10\,\% usa la heurística de costo mínimo con desempates
aleatorios. Esta siembra acelera la convergencia sin fijar la solución.

\textbf{Aptitud.} $f = 1/(Z + \lambda\cdot\mathrm{Pen})$, con $\lambda=10^6$ y $\mathrm{Pen}$ igual a
la suma de violaciones de oferta y de capacidad (plantas y acopios).

\textbf{Operadores.} Selección por torneo ($k=3$) o ruleta; cruce de uno o dos puntos con cortes
entre empresas (filas), lo que conserva la oferta de cada fila; mutación por transferencia de una
fracción aleatoria del flujo entre dos caminos de la misma empresa; elitismo de 2 individuos y parada
tras 50 generaciones sin mejora. Todos los parámetros se configuran desde la interfaz.

\textbf{Reparación.} Tras cada operador: (1)~se anulan flujos negativos o en rutas cerradas,
(2)~se reescala cada fila a su oferta y (3)~el exceso de cada recurso saturado (planta o acopio) se
mueve al camino más barato con cupo, empezando por las empresas que menos pierden. Así, casi todos los
individuos son factibles y la penalización solo actúa como respaldo.

\section{Arquitectura del sistema}
El sistema sigue una arquitectura por capas. El backend (FastAPI, Python~3.12) expone una API REST
con módulos independientes: \emph{red} (estructura y validación), \emph{exacto} (PuLP + CBC),
\emph{genético}, \emph{sensibilidad} (re-optimización paramétrica con HiGHS~\cite{highs}) y
\emph{criticidad} (ICE). Los datos persisten en PostgreSQL con migraciones versionadas (Alembic). El
frontend (React + TypeScript) ofrece pantallas de red, rutas, exacto, genético, sensibilidad, riesgo,
comparativo e historial, y un mapa real sobre OpenStreetMap~\cite{osm} con trazados viales de
OSRM~\cite{osrm}. Los reportes se exportan en PDF y Excel. Todo se ejecuta con
\texttt{docker compose up}; la integración continua corre el lint, las migraciones y las pruebas en
cada cambio.

\section{Datos de prueba regionales}
\textbf{Caso base.} 10 empresas (845~t/semana) y 4 plantas (950~t de capacidad) con distancias viales
estimadas y $k=350$~COP/t·km.

\textbf{Caso real del Meta.} 12 empresas en los municipios de mayor producción, con oferta estimada a
partir de la Evaluación Agropecuaria Municipal (EVA 2018)~\cite{eva} y del inventario bovino
ICA/FNG 2025~\cite{ica}, coordenadas DIVIPOLA~\cite{divipola} y distancias reales por carretera de
OSRM: 919{,}3~t/semana y 48 rutas.

\textbf{Centros de acopio.} Puerto López y San Martín (200~t, $h_t = 8000$~COP/t), sobre las vías
hacia Puerto Gaitán y Granada, con flete consolidado de 200~COP/t·km hacia las plantas.

\section{Resultados}
\subsection{Solución óptima}
En el caso base, $Z^* = 64\,691\,000$~COP/semana (CBC, 30~ms). Villavicencio y Restrepo quedan
saturadas, con precios sombra de $-17\,750$ y $-14\,250$~COP/t: cada tonelada adicional de capacidad
en Villavicencio ahorraría 17\,750~COP por semana. Granada (5~t libres) y Puerto Gaitán (55~t) tienen
precio sombra cero. En el caso real, $Z^* = 59\,295\,262$~COP/semana, sin planta ficticia.

\subsection{Análisis de sensibilidad}
Cada escenario se re-optimiza con ambos métodos (Tabla~\ref{tab:escenarios}). Un alza del 20\,\% en
el combustible eleva el costo en 7{,}3\,\%, sin cambiar la asignación. En cambio, un 10\,\% más de
oferta o una reducción del 30\,\% en Villavicencio superan la capacidad y obligan a usar la planta
ficticia: allí la capacidad, no el flete, es la restricción dominante. Los rangos de estabilidad se
calculan por re-optimización paramétrica (búsqueda y bisección), y la gráfica de tornado ordena los
parámetros por su impacto en $Z$.

\begin{table}[h]
\centering
\caption{Escenarios de sensibilidad en el caso base}
\label{tab:escenarios}
\begin{tabular}{@{}lrrr@{}}
\toprule
Escenario & $Z^*$ (MCOP) & $\Delta Z$ & Brecha AG \\
\midrule
Base & 64{,}69 & --- & 0{,}00\,\% \\
Combustible +20\,\% & 69{,}39 & +7{,}3\,\% & 0{,}00\,\% \\
Oferta +10\,\%$^{\dagger}$ & 98{,}35 & +52{,}0\,\% & 0{,}00\,\% \\
Villavicencio $-30$\,\%$^{\dagger}$ & 93{,}89 & +45{,}1\,\% & 0{,}00\,\% \\
\bottomrule
\multicolumn{4}{@{}l}{\footnotesize $^{\dagger}$ Usa la planta ficticia (penalización $10^6$ COP/t).}
\end{tabular}
\end{table}

\subsection{Comparativo exacto frente al AG}
En 10 corridas sobre el caso base (población 100, hasta 300 generaciones), el AG alcanzó el óptimo en
todas (brecha media $<10^{-7}$\,\%), con un tiempo medio de 1{,}9~s frente a 30~ms del exacto. En
instancias sintéticas reproducibles del Meta (Tabla~\ref{tab:escala}), la brecha crece con el tamaño y
el exacto sigue siendo más rápido: para un problema lineal de este tamaño, el método exacto domina. El
valor del AG está en que tolera extensiones no lineales (costos fijos, economías de escala) sin
cambiar su estructura.

\begin{table}[h]
\centering
\caption{Escalabilidad (semilla 42)}
\label{tab:escala}
\begin{tabular}{@{}lrrr@{}}
\toprule
Empresas $\times$ plantas & Exacto (ms) & AG (s) & Brecha AG \\
\midrule
$10\times4$ & 51 & 3{,}0 & 0{,}00\,\% \\
$30\times8$ & 61 & 4{,}0 & 0{,}19\,\% \\
$50\times10$ & 70 & 8{,}6 & 0{,}48\,\% \\
$100\times20$ & 216 & 9{,}8 & 0{,}91\,\% \\
\bottomrule
\end{tabular}
\end{table}

\subsection{Centros de acopio}
Con los dos acopios, $Z^*$ baja a 61\,761\,025~COP/semana (ahorro del 4{,}5\,\%): pasan 160~t por
Puerto López y 85~t por San Martín, ninguno saturado. El AG por caminos llega a una brecha de
0{,}64\,\% sin violaciones. Si los acopios no tienen cupo, el modelo vuelve exactamente a
64\,691\,000, lo que confirma que la generalización conserva el modelo original.

\subsection{Criticidad estructural (ICE)}
Se adaptó el índice de Vega-Linares, Beltrán-Tabares \emph{et al.}~\cite{ice},
$\mathrm{ICE}(i)=\alpha C_B(i)+\beta C_D^-(i)+\gamma(1-R(i))$, propuesto para dependencias entre
microservicios. En la red logística: $C_B$ es la fracción del flujo que recibe la planta (en un grafo
bipartito la intermediación clásica~\cite{freeman} es cero) o, para un acopio, la fracción de pares
empresa–planta cuyo camino más barato lo atraviesa; $C_D^-$ es la fracción de empresas que dependen
del nodo; y $R$ es la proporción de cinco mecanismos de mitigación: holgura $\ge 10$\,\%, precio
sombra cero, capacidad que puede bajar 10\,\% sin cambiar los duales, rutas alternativas a
$\le 10$\,\% de costo y respaldo por acopio. El fallo se simula cerrando el nodo y re-optimizando.

Con los pesos del artículo ($0{,}5$; $0{,}3$; $0{,}2$), el ranking del caso base es Villavicencio
(ICE $=0{,}49$; 240~t sin planta al cerrarla) $>$ Granada $>$ Restrepo $>$ Puerto Gaitán, con
$\rho$ de Spearman~\cite{spearman} de $0{,}80$ frente al impacto observado; en el caso real,
$\rho = 1{,}00$. Como en el artículo, el término $1-R$ reordena el ranking: Puerto Gaitán recibe casi el
mismo flujo que Restrepo, pero su holgura y su precio sombra cero lo hacen menos crítico. Una búsqueda
en malla de los pesos muestra que, en el caso base, la dependencia $C_D^-$ sola explica mejor el
impacto ($\rho=0{,}95$). Con solo 4 a 6 nodos, estos resultados son evidencia preliminar y no una
validación concluyente; constituyen la prueba empírica que el artículo dejó como trabajo futuro.

\section{Conclusiones y trabajo futuro}
El modelo exacto resuelve el problema en milisegundos y entrega la información de gestión más útil:
qué plantas están saturadas y cuánto vale ampliarlas. El AG, con reparación por recursos, alcanza el
óptimo en el caso base y se mantiene por debajo del 1\,\% de brecha en instancias de 100 empresas. La
red es sensible sobre todo a la capacidad: la ampliación de Villavicencio es la inversión prioritaria,
y los centros de acopio con flete consolidado reducen el costo un 4{,}5\,\%. El ICE identifica la
misma planta crítica que la simulación de fallas y ofrece una lectura de riesgo comprensible para la
gerencia.

Como trabajo futuro se propone calibrar el ICE con incidentes reales de la empresa, alimentar con
estas características una red neuronal de grafos (GAT), como plantea~\cite{ice}, e incorporar costos
fijos de apertura de acopios, que convierten el problema en uno entero mixto donde el AG gana ventaja.

\section*{Declaración de uso de inteligencia artificial}
% TODO(equipo): ajustar esta declaración al uso real.
El equipo usó asistentes de inteligencia artificial como apoyo en la programación, las pruebas y la
redacción. Todas las formulaciones, resultados y conclusiones fueron revisados y verificados por los
autores, que asumen la responsabilidad del contenido.

\begin{thebibliography}{99}
\bibitem{hillier} F.~S. Hillier y G.~J. Lieberman, \emph{Introduction to Operations Research}, 10th ed.
New York, NY, USA: McGraw-Hill, 2015.
\bibitem{goldberg} D.~E. Goldberg, \emph{Genetic Algorithms in Search, Optimization, and Machine
Learning}. Reading, MA, USA: Addison-Wesley, 1989.
\bibitem{pulp} S.~Mitchell, M.~O'Sullivan e I.~Dunning, ``PuLP: A linear programming toolkit for
Python,'' Univ. of Auckland, Auckland, New Zealand, 2011.
\bibitem{cbc} J.~Forrest y R.~Lougee-Heimer, ``CBC user guide,'' en \emph{Emerging Theory, Methods, and
Applications}, INFORMS, 2005, pp.~257--277.
\bibitem{highs} Q.~Huangfu y J.~A.~J. Hall, ``Parallelizing the dual revised simplex method,''
\emph{Math. Program. Comput.}, vol.~10, no.~1, pp.~119--142, 2018.
\bibitem{ice} L.~A. Vega-Linares, D.~F. Beltrán-Tabares, O.~L. Vega-Márquez y N.~E. Suat-Rojas,
``Redes neuronales de grafos para resiliencia en microservicios,'' en \emph{Memorias del VI Congreso
Internacional de Ciencias Básicas e Ingeniería -- CICI 2026}, Rev. Punto de Inflexión, ISSN 3028-6425,
Universidad de los Llanos, Villavicencio, 2026.
\bibitem{freeman} L.~C. Freeman, ``A set of measures of centrality based on betweenness,''
\emph{Sociometry}, vol.~40, no.~1, pp.~35--41, 1977.
\bibitem{spearman} C.~Spearman, ``The proof and measurement of association between two things,''
\emph{Amer. J. Psychol.}, vol.~15, no.~1, pp.~72--101, 1904.
\bibitem{eva} Ministerio de Agricultura y Desarrollo Rural y UPRA, ``Evaluaciones Agropecuarias
Municipales (EVA),'' 2018. [En línea]. Disponible: datos.gov.co % TODO(equipo): URL exacta
\bibitem{ica} Instituto Colombiano Agropecuario y FNG, ``Censo pecuario nacional,'' 2025. [En línea].
Disponible: ica.gov.co % TODO(equipo): URL exacta
\bibitem{divipola} DANE, ``Codificación de la División Político-Administrativa de Colombia
(DIVIPOLA),'' 2024. [En línea]. Disponible: dane.gov.co % TODO(equipo): URL exacta
\bibitem{osrm} D.~Luxen y C.~Vetter, ``Real-time routing with OpenStreetMap data,'' en \emph{Proc.
19th ACM SIGSPATIAL Int. Conf. Adv. Geographic Inf. Syst.}, 2011, pp.~513--516.
\bibitem{osm} OpenStreetMap contributors, ``OpenStreetMap,'' 2026. [En línea]. Disponible:
\url{https://www.openstreetmap.org}
\end{thebibliography}

\end{document}
